\documentclass[aspectratio=169,10pt,t]{beamer}
\usepackage{ac1}
\aulainfo{08}{Iteração}{Laços while e for, laços contados e condições complexas}

\begin{document}
\capa

\begin{frame}{Roteiro}
  \begin{itemize}
    \item Por que repetir? Iteração (Cap.~3.1.3 e 3.5)
    \item O comando \texttt{while}
    \item Laços contados com \texttt{for} e \texttt{range} (3.6)
    \item Padrões: acumulador, contador, sentinela
    \item Laços encaixados
    \item Condições complexas e desafios de lógica
  \end{itemize}
\end{frame}

\begin{frame}{Iteração: repetir um bloco}
  \begin{columns}[T]
    \column{0.55\textwidth}
    \begin{itemize}
      \item Computadores são ótimos em \textbf{repetir} tarefas sem se cansar
      \item Somar \(1+2+\cdots+n\), processar mil medições, tabular uma função
      \item Escrever 1000 \texttt{print}s é inviável; um laço resolve em 3 linhas
    \end{itemize}
    \column{0.42\textwidth}
    \begin{block}{Dois tipos de laço}
      \begin{itemize}\small
        \item \texttt{while}: repete \emph{enquanto} uma condição for verdadeira
        \item \texttt{for}: repete para cada elemento de uma sequência
      \end{itemize}
    \end{block}
  \end{columns}
\end{frame}

\begin{frame}{Fluxograma do \texttt{while}}
  \begin{center}
  \begin{tikzpicture}[scale=0.9,transform shape,node distance=0.9cm,
      proc/.style={draw=uerjazul,fill=uerjcinza,minimum width=2.8cm,minimum height=0.7cm,font=\small},
      dec/.style={draw=uerjazul,diamond,aspect=2.5,fill=uerjcinza,font=\small},
      seta/.style={-{Stealth},thick,uerjazul}]
    \node[dec] (c) {condição?};
    \node[proc,right=2.6cm of c] (b) {bloco do laço};
    \node[proc,below=1.2cm of c] (f) {continua};
    \draw[seta] ([yshift=0.8cm]c.north)--(c);
    \draw[seta] (c)--node[above,font=\scriptsize]{verdadeira}(b);
    \draw[seta] (b)|-([yshift=0.5cm]c.north);
    \draw[seta] (c)--node[left,font=\scriptsize]{falsa}(f);
  \end{tikzpicture}
  \end{center}
  \begin{alertblock}{Cuidado com o laço infinito}
    Algo dentro do bloco precisa \emph{modificar} a condição; caso contrário o programa nunca termina (\texttt{Ctrl+C} interrompe).
  \end{alertblock}
\end{frame}

\begin{frame}[fragile]{\texttt{while}: soma de 1 até \(n\)}
\begin{lstlisting}
n = int(input('Qual o número final? '))
soma = 0
i = 1
while i <= n:
    soma = soma + i
    i = i + 1
print('A soma de 1 até ', n, ' = ', soma)
\end{lstlisting}
  \begin{itemize}
    \item \texttt{soma}: \alert{acumulador} (começa em 0)
    \item \texttt{i}: \alert{contador} (começa em 1 e cresce até ultrapassar \(n\))
    \item Teste na mão (\(n=4\)): \(i=1,\,soma=1\); \(i=2,\,soma=3\); \(i=3,\,soma=6\); \(i=4,\,soma=10\)
  \end{itemize}
  \fonte{code/Cap3/prog.3.11\_soma\_numeros.py (Exercício 3.11)}
\end{frame}

\begin{frame}[fragile]{Validando a entrada}
\begin{lstlisting}
n = int(input('Qual o número final? '))
if n > 0:
    soma = 0
    i = 1
    while i <= n:
        soma = soma + i
        i = i + 1
    print('A soma de 1 até ', n, ' = ', soma)
else:
    print('Use números positivos')
\end{lstlisting}
  \begin{itemize}
    \item Combina \texttt{if} (validação) com \texttt{while} (cálculo)
    \item Aviso: \texttt{n} negativo ou zero não deve gerar resultado
  \end{itemize}
  \fonte{code/Cap3/prog.3.12\_soma\_numeros\_v2.py}
\end{frame}

\begin{frame}[fragile]{Sentinela: repetir até o usuário mandar parar}
\begin{lstlisting}
n = int(input('Qual o número final? (0 para terminar) '))
while n != 0:
    if n > 0:
        soma = 0
        i = 1
        while i <= n:
            soma = soma + i
            i = i + 1
        print('A soma de 1 até ', n, ' = ', soma)
    else:
        print('Use números positivos')
    n = int(input('Qual o número final? (0 para terminar) '))
\end{lstlisting}
  \small O valor \texttt{0} é o \alert{sentinela}. Note os \textbf{laços encaixados} e a releitura de \texttt{n} no fim.\\
  \fonte{code/Cap3/prog.3.13\_soma\_numeros\_v3.py}
\end{frame}

\begin{frame}[fragile]{Média de números digitados (Exercício 3.12)}
\begin{lstlisting}
soma = 0
quantidade = 0
valor = float(input("Digite um valor (negativo para parar): "))
while valor >= 0:
    soma += valor
    quantidade += 1
    valor = float(input("Digite um valor (negativo para parar): "))
if quantidade > 0:
    print("Média:", soma / quantidade)
else:
    print("Nenhum valor digitado")
\end{lstlisting}
  \begin{itemize}
    \item \texttt{soma += valor} é o mesmo que \texttt{soma = soma + valor}
    \item Proteja a divisão: se \texttt{quantidade} for 0, evita \texttt{ZeroDivisionError}
  \end{itemize}
\end{frame}

\begin{frame}[fragile]{Laços contados: \texttt{for} e \texttt{range}}
\begin{lstlisting}
for i in ["Vermelho", "Verde", "Azul"]:
    print(i)
for i in range(1, 10):
    print(i, end=' ')
\end{lstlisting}
\begin{lstlisting}[style=saida]
Vermelho
Verde
Azul
1 2 3 4 5 6 7 8 9
\end{lstlisting}
  \begin{center}\small
  \begin{tabular}{ll}
    \toprule
    \texttt{range(5)} & 0, 1, 2, 3, 4\\
    \texttt{range(1, 10)} & 1, 2, \ldots, 9 (o final \alert{não} entra)\\
    \texttt{range(0, 361, 30)} & 0, 30, 60, \ldots, 360\\
    \texttt{range(10, 0, -2)} & 10, 8, 6, 4, 2\\
    \bottomrule
  \end{tabular}
  \end{center}
  \fonte{code/Cap3/prog.3.14\_cores.py, prog.3.15\_contagem.py}
\end{frame}

\begin{frame}[fragile]{Fatorial com \texttt{for}}
\begin{lstlisting}
n = int(input('Qual o número para calcular o fatorial? '))
fat = 1
for i in range(1, n+1):
    fat = fat * i
print('O fatorial de', n, 'é', fat)
\end{lstlisting}
  \begin{itemize}
    \item \(n! = 1\times2\times\cdots\times n\); acumulador multiplicativo começa em \alert{1}, não em 0
    \item \texttt{range(1, n+1)}: o limite superior é excluído
    \item Compare com a soma por \texttt{while}: quando se sabe quantas vezes repetir, \texttt{for} é mais simples (Exercício 3.13)
  \end{itemize}
  \fonte{code/Cap3/prog.3.17\_fatorial.py}
\end{frame}

\begin{frame}[fragile]{Quando usar \texttt{for} ou \texttt{while}?}
  \begin{columns}[T]
    \column{0.48\textwidth}
    \begin{block}{\texttt{for}}
      Número de repetições conhecido ou percorrer uma sequência.\\[0.5ex]
      \footnotesize Tabela de 0 a 360 graus; notas de uma turma.
    \end{block}
    \column{0.48\textwidth}
    \begin{block}{\texttt{while}}
      Repetir até que algo aconteça.\\[0.5ex]
      \footnotesize Ler até valor válido; aproximar até o erro ficar pequeno.
    \end{block}
  \end{columns}
  \vspace{1.5ex}
  \texttt{break} sai do laço; \texttt{continue} pula para a próxima volta:
\begin{lstlisting}
for n in range(1, 100):
    if n % 7 == 0 and n % 5 == 0:
        print("Primeiro múltiplo de 35:", n)
        break
\end{lstlisting}
\end{frame}

\begin{frame}[fragile]{Laços encaixados: padrões}
\begin{lstlisting}
for linha in range(1, 5):
    for coluna in range(linha):
        print('*', end='')
    print()
\end{lstlisting}
\begin{lstlisting}[style=saida]
*
**
***
****
\end{lstlisting}
  \begin{itemize}
    \item O laço interno \emph{completo} executa a cada volta do externo
    \item Exercício 3.15: gerar outros padrões com laços encaixados
  \end{itemize}
\end{frame}

\begin{frame}{Condições complexas: desafios de lógica}
  \begin{exampleblock}{Desafio 1}
    Mostre todos os anos \textbf{bissextos} entre 1900 e 2030. (Revisite a condição da Aula 7.)
  \end{exampleblock}
  \begin{exampleblock}{Desafio 2}
    Mostre os números de 1 a 100 que são divisíveis por 3 \textbf{ou} 5, mas \textbf{não} por ambos.
  \end{exampleblock}
  \begin{exampleblock}{Desafio 3}
    Leia números até digitar 0 e informe o maior e o menor.
  \end{exampleblock}
  \vspace{1ex}
  \small Estratégia: escreva primeiro a \emph{condição} em português, depois traduza para \texttt{and}/\texttt{or}/\texttt{not}, e teste os extremos.
\end{frame}

\begin{frame}[fragile]{Desafio 1: bissextos em um intervalo}
\begin{lstlisting}
for ano in range(1900, 2031):
    if (ano % 4 == 0 and ano % 100 != 0) or ano % 400 == 0:
        print(ano, end=' ')
\end{lstlisting}
\begin{lstlisting}[style=saida]
1904 1908 1912 ... 1996 2000 2004 ... 2028
\end{lstlisting}
  \begin{itemize}
    \item 1900 \alert{não} aparece (divisível por 100, não por 400); 2000 aparece
    \item Dica: conte quantos anos foram impressos usando um contador
  \end{itemize}
\end{frame}

\exercicios{3}{%
  \begin{itemize}\small
    \item \texttt{while}: \textbf{3.11, 3.12}
    \item \texttt{for}: \textbf{3.13} (refazer 3.11), \textbf{3.14} (tabela Celsius \(\to\) Fahrenheit), \textbf{3.15} (padrões com laços encaixados)
  \end{itemize}}

\resumo{
  \item \texttt{while} repete enquanto a condição for verdadeira: cuidado com laços infinitos
  \item \texttt{for} com \texttt{range} é ideal para laços contados
  \item Padrões: contador, acumulador, sentinela
  \item Laços podem ser encaixados; \texttt{break} e \texttt{continue} controlam o fluxo
  \item Próxima aula (prática): laboratório de problemas reais
}

\end{document}
